\chapter{連結グラフ}
    \section{諸定理}
        \subsection{\texorpdfstring{$n$}{n}頂点の連結グラフは少なくとも\texorpdfstring{$n-1$}{n-1}本の辺をもつ}
            \begin{proof}
            \quad\par
                $n$に関する帰納法で示す。$n=2$の時は明らか。$n-1$まで成り立つと仮定して$n$のとき成り立つことを示す。
                \par
                グラフから辺を任意に1本ずつ取り去ってゆき、$c$本目で初めて非連結になったとする。このときグラフは2つの連結成分$A,B$に分かれる(3つ以上に分かれることはない。3つ以上の連結成分を1本の辺で橋渡しするなど不可能である)。両者の点数をそれぞれ$a,n-a\;(1\leq a\leq n-1)$とすると帰納法の仮定より$A$は$a-1$本以上の、$B$は$n-a-1$本以上の辺をそれぞれもつ。よって元々のグラフは$(a-1)+(n-a-1)+c = n-2+c \geq n-1$本以上の辺をもっていたことになる。
            \end{proof}
        \subsection{任意の2つの頂点間の最短距離が2である無向グラフは少なくとも\texorpdfstring{$n-1$}{n-1}本の辺をもつ}
            \begin{proof}
                $n=2,3$の場合は明らかである。以下$n\geq4$とする。
                無向グラフ$G$が$n$個の点をもち、任意の頂点間の最短距離が2であり、$G$から1つでも辺を取り去ると、ある2点間の最短距離が3以上になってしまうとする。
                $G$から任意の点を1つ選んで$v_0$とする。
                $v_0$に隣接しているから成る集合を$S_0$とする。
                $S_0$に属していない点のうち、「$S_0$中にある点が存在して、それと隣接している」全ての点について、その点から$S_0$に繋がる辺1本の接続先を$v_0$に変更する。
                それらの点を$S_0$に加えた集合を$S_1$とする。
                以降同様にして$S_2,S_3,\dotsc,...$としてゆくと、$v_0$以外の全ての点が$v_0$に隣接する。
                そして$v_0$以外のいくつかの点同士の間に辺が存在し得る。
                以上の操作は$G$の辺の数を保つ。
                $v_0$とそれ以外の点を結ぶ辺のみを残して他の辺を除去したグラフは、任意の2点間の最短距離が2である。
                以上より、定理の主張が成り立つ。
            \end{proof}
        \subsection{\texorpdfstring{$n$}{n}頂点グラフが\texorpdfstring{$\frac{1}{2}(n-1)(n-2)$}{(n-1)(n-2)/2}本より多くの辺を持つなら連結である}
            \begin{proof}
                \quad\par
                $n$に関する帰納法で示す。$n=1,2$のときは明らかに成り立つ。$n=m-1\;(m\geq3)$のとき成り立つものとして$n=m$のとき成り立つことを示す。
                \par
                $n=m$のときのグラフを$G$とする。辺の数$k$について$k>\frac{1}{2}(m-1)(m-2)$という前提のもとで仮に$G$が非連結であるとする。$k>0$なので次数が1以上の点が必ず存在するので任意に1つ選んで$v$とし、$G$から$v$とそれに接続する辺を取り去ったグラフを$G'$とするとこれは非連結である($\cdots$\textcircled{1})。(※$\deg(v)\geq1$であるところがミソ。$\deg(v)=0$だと、もし不運なことに$v$が唯一の孤立点だったら$G'$が連結になってしまう。)$G'$は非連結であるから$\deg(v)\leq m-2$である(もし$\deg(v)=m-1$なら$v$は他の全ての頂点と隣接していることになり、$G$が連結になってしまう)。$G'$の辺の数は$k'=k-\deg(v)$であり、
                \[k' = k-\deg(v) > \frac{1}{2}(m-1)(m-2) - \deg(v) > \frac{1}{2}(m-1)(m-2) - (m-2)\]
                ここで
                \[\left[\frac{1}{2}(m-1)(m-2) - (m-2)\right] - \frac{1}{2}[(m-1)-1][(m-1)-2] = m-2 - (m-2) = 0\]
                であるから
                \[k' > \frac{1}{2}[(m-1)-1][(m-1)-2]\]
                これと帰納法の仮定から$G'$は連結でなくてはならないが、このことは\textcircled{1}と矛盾する。よって$G$が非連結であるという仮定は否定される。
            \end{proof}
        \subsection{Oreの定理}
            \begin{shadebox}
                Oreの定理\\
                $n(\geq3)$点以上からなる単純グラフ$G$の隣接していない任意の2点$v,w$について$\deg(v)+\deg(w)\geq n$であれば$G$はハミルトングラフである。
            \end{shadebox}
            \begin{proof}
                \quad\par
                まず、$G$が連結であることを示そう。仮に$G$が2つの連結成分から成るとする。一方($A$とする)の点数を$a$とすれば他方($B$とする)の点数は$n-a$である。$A,B$から任意の点を1つずつ選んで$v,w$とすると$\deg(v)\leq a-1,\;\;\deg(w)\leq n-a-1$であるから$\deg(v)+\deg(w)\leq n-2$となって定理の仮定に反する。このように$G$が2つの連結成分から成る場合でさえ許されないのだから、3つ以上など以ての外である。よって$G$は連結である。
                \par
                次に$G$がハミルトングラフであることを示す。$G$の単純で最長な道を$p : (u_0,u_1,\dotsc,u_{k-1},u_k)$とする。この中の点は全て異なるので$k+1\leq n$である。この道は最長だから$u_0$および$u_k$の隣接接点は全てこの道の中にある。集合$A:\{0\leq i\leq k-1|u_{i+1}\text{は}u_0\text{の隣接接点}\}$、$B:\{0\leq i\leq k-1|u_i\text{は}u_k\text{の隣接接点}\}$を考えると、当然$n(A\cup B)\leq k\leq n-1$である。然るに定理の仮定から$n(A)+n(B)\geq n$であるので$A$と$B$は共通部分を持つ,すなわち、ある$i\;(0\leq i\leq k-1)$が存在して$u_i$は$u_k$の隣接接点で、$u_{i+1}$は$u_0$の隣接接点である。
                \begin{figure}[H]
                    \centering
                    \includegraphics[keepaspectratio, scale=0.7]
                    {\currfiledir/figs/OreTheorem/OreTheorem.pdf}
                    \caption{様子}
                \end{figure}
                ここで閉路$c : (u_0,u_{i+1})+p[u_{i+1},u_k]+(u_k,u_i)+p[u_i,u_0]$を考える。但し$p[u,v]$は道$p$の$u$から$v$への部分を表す。実はこの閉路$c$がハミルトン閉路になっている。
                \par
                $c$がハミルトン閉路であることを示そう。仮にこれがハミルトン閉路でないとすると、$G$上のある点$v$が存在してこれは$c$に含まれない。先述の如く$G$は連結なので$v$はどうにかして$c$と繋がっている,すなわち、これまたある点$w$が存在して(※たまたま$v$と一致することもある)これは$c$に属さないが$c$上のある点$u_t$に隣接している。このとき、$(w,u_t)$から始めて$c$を上手く回って$u_t$の手前まで行く単純な道を考えると(※簡単に考えられる。例えば上の図で$w$が$u_2$に隣接している場合を考えてみればいい。)これは$p$より長くなってしまうので、$p$が最長であったという前提に矛盾する。よって$G$上の点で$c$に含まれないものは無かったということになる。以上より$c$はハミルトン閉路である。
            \end{proof}